\section{DSU и остовные деревья}

\subsection{Домашнее задание}
\begin{enumerate}
  \item
    Дан неориентированный граф с весами на вершинах, в котором в начальный момент нет ребер.
    Вершин $V$ штук. Нужно обрабатывать следующие запросы:
    \begin{itemize}
      \item соединить ребром пару вершин.
      \item по вершине узнать размер связной компоненты, которой вершина принадлежит, и самую
            легкую вершину в этой компоненте.
    \end{itemize}
    $\O(\log^* V)$ на запрос в среднем, \texttt{online}.

  \item
    Есть отрезок длины $n$ и $m$ запросов вида $paint(l, r, c)$ -- перекрасить подотрезок $[l..r]$ в цвет $c$ (все координаты целые).
    Определите состояние отрезка после выполнения всех запросов за $\O((n + m)\log^*n)$.

  \item
    За $O(E \log V)$ для каждого ребра определите, существует ли минимальный остов, содержащий его.

  \item
    Вам нужно передать с одного компьютера на другой $n$ файлов. Каждый файл -- битовый вектор
    размера $m$. Файл можно передать либо просто как последовательность битов, либо как $diff$
    с другим файлом, уже пересланным ранее. В первом случае надо будет передать $m$ бит, во
    втором - $A+B*d$ бит, где $d$ - число различающихся битов, $A$ и $B$ -- некоторые
    константы. Найти минимальное число битов, которое нужно передать, чтобы отправить все
    файлы с одного компьютера на другой. $\O(mn^2)$.

\end{enumerate}

\newpage
\subsection{Решение}

\begin{enumerate}

\item

    Для решения задачи воспользуемся СНМ.
    Каждая компонента связности представляет из себя отдельное множество, где вся информация об этой компоненте хранится в корневой вершине.
    Для каждой вершины храним информацию о размере связной компоненты, которой она принадлежит, и самой легкой вершине в этой компоненте.

    Пусть для вершины $v$:
    \begin{itemize}
        \item $p[v], r_v$ -- корневой родитель вершины $v$.
        \item $size[v]$ -- размер компоненты, которой принадлежит вершина $v$.
        \item $min\_weight[v]$ -- минимальный вес вершины, подсчитанный в какой-то момент времени.
    \end{itemize}

    При запросе на соединение вершин $a$ и $b$ найдем их корневых родителей $r_a$ и $r_b$ применяя при этом сжатие пути.

    Если вершины принадлежат одной компоненте ($r_a = r_b$), то ничего не делаем.
    Иначе подцепим более легкую вершину к более тяжелой, обновим информацию в новом корне:
    \begin{itemize}
        \item $size[r_n] = size[r_a] + size[r_b]$.
        \item $min\_weight[r_n] = \min(min\_weight[r_a], min\_weight[r_b])$.
    \end{itemize}

    При запросе на информацию о вершине $v$ вернем информацию из корневой вершины $r_v$.

    Благодаря, сжатию пути и объединению по размеру, мы можем обрабатывать запросы за $\O(\log^* V)$ в среднем.

\item

    Для решения задачи создадим систему непересекающихся множеств, где каждая клетка отрезка будет хранить ссылку на
    ближайшую неокрашенную клетку справа. Таким образом, изначально каждая клетка указывает на саму себя, а после
    покраски первого подотрезка будет указывать на первую клетку после конца подотрезка.

    Теперь для корректного решения задачи будем рассматривать запросы в обратном порядке, чтобы
    задача перекрашивания подотрезка сводилась к покраске пустых клеток в цвет $c$ на отрезке $[l, r]$.
    Таким образом, во время обработки распросов мы красим каждую клетку ровно один раз и перепрыгиваем через уже покрашенные.
    
    

\item

    Для задачи определения, содержится ли каждое ребро в минимальном остовном дереве, построим минимальное остовное дерево с использованием алгоритма Прима.
    Все его ребра гарантированно будут входить в результат, однако для проверки принадлежности ребра к любому минимальному остову этого недостаточно.
    Поскольку у графа может быть несколько минимальных остовных деревьев. Попробуем за счет каждого ребра вне остова нарушить критерий уникальности
    минимального остова.

    Пусть $e = (u, v)$ -- ребро, которое не содержится в построенном минимальном остове.
    Рассмотрим максимальное ребро на пути от $u$ до $v$ внутри остова $e' = (u', v')$:

    \begin{itemize}
        \item Если $w(e) > w(e')$, сценарий невозможен, иначе ребро было бы частью найденного остова.
        \item Если $w(e) < w(e')$, сценарий увеличит вес остова, что противоречит минимальности.
        \item Если $w(e) = w(e')$, мы можем заменить ребро $e'$ на $e$ и получить новый минимальный остов.
    \end{itemize}

    Для поиска максимального ребра на пути используем heavy-light декомпозиции за $\O(\log V)$ на запрос.
    Общее выполнение для всех ребер займет $\O(E \log V)$.

    Таким образом, построение минимального остова работает за $\O(E \log V)$, построение heavy-light декомпозиции
    за $\O(E)$ с последующими запросами за $\O(\log V)$. Минимальный остов и heavy-light декомпозиция строятся только один раз.

\item

    Для решения задачи можно представить файлы как вершины графа, а различия между ними — как ребра.
    Вес ребра между вершинами \(u\) и \(v\) определяется как \(w(u, v) = A + B \cdot d(u, v)\),
    где \(d(u, v)\) — это количество различных битов между файлами \(u\) и \(v\).

    Однако такой граф не учитывает, что файлы можно передавать в исходном формате и что разницу можно передавать только с уже переданными файлами.
    Для этого введем виртуальный корень \(r\) и добавим к нему ребра ко всем вершинам с весом, равным размеру файла, если передавать его в исходном формате.

    Задача сводится к поиску минимального остовного дерева в графе, начинающемся с виртуального корня \(r\).
    Для этого можно использовать алгоритм Прима и после построения посчитать общее количество бит, основываясь на весах ребер в остове.

    Общая сложность алгоритма составляет \(\O(mn^2)\) на построение графа, а для поиска минимального остовного
    дерева используется наивный алгоритм Прима с сложностью \(\O(n^2)\), так как у нас плотный граф.

\end{enumerate}
